\subsection{Open questions}
%==========================================================================

\paragraph{Well-posedness.} Three quantities in the formulation are not measured and are
not physiological constants: $\mtol$ (clinical convention), $\spk$ (varies with the
precipitating event), and the map from intracellular substrate to plasma biomarker. $\thc$
is therefore an operational definition; it is not a physical parameter. Whether the
decision built on it is invariant to that arbitrariness is the central question.

\paragraph{From a threshold to a hazard.}
\label{sec:hazard} This is the reformulation referred to in
Remark~\ref{rem:ceiling} and Remark~\ref{rem:intrinsic_limit}, and it is
not attempted here. Proposition~\ref{prop:intrinsic} removes $\mtol$ from the definition
of $\thc$ but leaves $\spk$, and no sharper deterministic threshold will remove it, because
the quantity actually wanted is a rate. Treat the precipitating events as a marked
point process with arrival rate $\nu$ and amplitudes drawn from a distribution
\emph{estimated from the records rather than postulated}, and ask for the rate at which the
metabolite crosses the separatrix:
\begin{equation}
\lambda_{\mathrm{MDE}}(\theta) \;=\; \nu \cdot
\mathbb{P}\!\left(A > A_{\mathrm{crit}}(\theta) \mid \theta \right),
\label{eq:hazard}
\end{equation}
where $A$ is the amplitude of an episode and $A_{\mathrm{crit}}(\theta)$ is the smallest
excursion in the load that carries the patient from the attractor across the separatrix.
The equilibrium reduction already supplies $A_{\mathrm{crit}}$: it is the margin of
Figure~\ref{fig:C5},
converted from metabolite into load.

The construction is partly pre-empted. The natural history study already fits a hazard of
exactly this
shape: a generalised additive \emph{mixed} model on 3-HP concentration, with a smooth term
in the days since the previous event and a participant-level random
effect~\cite{natural_history}. % C28
What is proposed here is a substitution within that machinery, of a mechanistic $\theta$
with $A_{\mathrm{crit}}(\theta)$ supplied by the model, for a measured biomarker as the
covariate carrying the risk. The statistical apparatus already exists; what it is
regressed on is the open question.

Four consequences:

\begin{enumerate}[leftmargin=*, itemsep=2pt]
\item $\spk$ leaves the definition of the threshold. A single assumed worst case is
      replaced by a distribution, so nothing has to be asserted about the severity of the
      worst illness a given patient will meet.
\item The endpoint becomes the reliable observable. Dated admissions are coded
      consistently, whereas the metabolite measurements are sparse and stand in an assumed
      relation to the enzyme's substrate.
\item The recurrent-event structure is used. Patients do not have one
      event each: $49$ of $50$ participants in the natural history cohort reported at least
      one, at a mean annualised frequency of $1.0$~\cite{natural_history}, % C08
      and the rate varies enough between them that the study carries a participant-level
      random effect. A hazard formulation uses every event and every event-free interval a
      patient contributes, and makes the spread between patients a parameter to be estimated
      instead of noise to be averaged over.
\item The count model becomes \emph{derived}. A Poisson or
      negative-binomial model imposed on admission counts is a modelling choice with
      nothing behind it; under \eqref{eq:hazard} the counting process follows from the
      arrival process and the crossing probability, and its dispersion has a mechanism.
\end{enumerate}

It would also close a gap the deterministic threshold leaves open in either structure. Both
return a \emph{condition}, written in quantities nobody measures, and the records contain a
\emph{frequency}, which is the one thing they hold reliably. $\lambda_{\mathrm{MDE}}(\theta)$
is a frequency, and it is the unit the two can be compared in.

\paragraph{What the marker choice still leaves open.} Section~\ref{sec:marker_orig} now fixes the
marker at 3-HP, on the endpoint association and on the sign of the perturbation, which is
what a decision-framed argument can settle from published data. It does not settle what the
identifiability analysis would want to compare the candidates on, and that gap is worth
stating because it is the one a record could close.

The residual capacity enters through the intercept, which by Remark~\ref{rem:fibre} locates
it only along a fibre while fixing its relative precision, and by \eqref{eq:floor} that
precision is governed by $n$ and by $\bar m/\Delta$, the level a patient sits at relative to
the excursion they make about it. Both are set by how the patient is managed, on that
marker's own scale, so the object to be compared is a marker \emph{together with the
management that holds a patient in a band on it}, and not a marker on its own. The
candidates plausibly differ there, in analytical variability, in how far they move between
visits, and in how nearly affine the clearance curve is over the range each visits. None of
that is established here, and assembling it against the same records the rest of this note
relies on is a well-posed piece of work.

Two things sharpen it rather than defer it. Remark~\ref{rem:panel} says the marker panel is
worth buying only after the load record is, which orders the two pieces of work.
Section~\ref{sec:discriminant} adds the gap that no marker is currently known to track the
state through a decompensation at all, so the marker fixed here is fixed for the stable
patient and the excursion, not for the crisis.

\paragraph{Identification.} Under what conditions is a parameter identifiable when the
excitation required to identify it is the regime that a safety constraint forbids? The
excitation is supplied by the schedule of the selective perturbation itself, but the
schedule must be selected \emph{before} any response data exist: the information arrives
after the moment at which it is required.

\paragraph{Design.} The safety constraint is active only at the trough preceding each dose
of the perturbation, where the total clearance capacity is at its lowest and the margin
$\theta - \thc$ is smallest. Designing a sequence of administration times that guarantees
this for every parameter value in an uncertainty set, while simultaneously contracting that
set, is dual control with a state constraint whose defining parameter is the unknown.
Violation is not an off-specification excursion but a hospitalisation carrying a risk of
permanent neurological cost, so the appropriate formulation is a chance constraint rather
than a quadratic objective.

\paragraph{Confounding of the accessible input.}
\label{sec:confounding}
The perturbation actually present in the
records is carnitine, and it is endogenous: the dose responds to the patient's state. The
design and identification questions above are therefore posed on an input that no protocol
fixed in advance. Whether the quasi-exogenous variation listed in
Section~\ref{sec:instruments}, namely weight-based recalculation after growth, supply
interruptions and changes of formulation, is sufficient to support the argument is a
question for the records, and it is answered patient by patient.
The load route is asked the same question and costs more: a change of load carries the same
determinant as a channel perturbation, and it needs both loads known where the channel
route needs one.

\paragraph{Reduction.} A model reduction preserving the input--output response to periodic
impulsive inputs does not obviously yield a reduced model whose parameters are identifiable
from sparse sampling. These are distinct objectives and may conflict.

\paragraph{Multi-parameter case.} The scalar analysis above is complete. With several
unknowns entering simultaneously, $\theta$, the population parameters, $\lambda$, and one
amplitude per documented illness episode, attractivity of the estimator manifold becomes a
matrix condition rather than a sign condition, and is not addressed here.

\paragraph{The therapy, when it arrives.} The analysis of this note is carried out on
carnitine because that is the selective input present in the accessible cohort. By
Proposition~\ref{prop:channel} nothing in the identification problem distinguishes the two
channels: the determinant depends on the metabolite levels before and after the
perturbation and not on which channel produced them. What the mRNA therapy would add is
therefore not new structure but a better \emph{statistical} position. Its schedule is fixed
by protocol, with no reference to the patient's state, so it supplies a genuinely exogenous
selective input, on the enzymatic channel, in a cohort where the only such input at present
is endogenous. The instrumental apparatus needed to defend a carnitine-based argument would
become unnecessary. The conditioning, which is the actual obstruction, is unchanged.
